\documentclass[11pt]{article}

% Change "review" to "final" to generate the final (sometimes called camera-ready) version.
% Change to "preprint" to generate a non-anonymous version with page numbers.
\usepackage[]{acl}
\usepackage{times}
\usepackage{latexsym}

% For proper rendering and hyphenation of words containing Latin characters (including in bib files)
\usepackage[T1]{fontenc}
% For Vietnamese characters
% \usepackage[T5]{fontenc}
% See https://www.latex-project.org/help/documentation/encguide.pdf for other character sets

% This assumes your files are encoded as UTF8
\usepackage[utf8]{inputenc}
\usepackage{microtype}
\usepackage{inconsolata}
\usepackage{graphicx}
\usepackage{subfiles}
\usepackage{float}

% If the title and author information does not fit in the area allocated, uncomment the following
%
%\setlength\titlebox{<dim>}
%
% and set <dim> to something 5cm or larger.

\title{Group 47 Progress Report:\\Single Class Movie Genre Classifier}


\author{Nour Fawaz, Daniel Maurer, Lily Porter, Sriya Dhanvi Mokhasunavisu\\
  \texttt{\{fawazn,maured2,portel7,mokhasus\}@mcmaster.ca} }

%\author{
%  \textbf{First Author\textsuperscript{1}},
%  \textbf{Second Author\textsuperscript{1,2}},
%  \textbf{Third T. Author\textsuperscript{1}},
%  \textbf{Fourth Author\textsuperscript{1}},
%\\
%  \textbf{Fifth Author\textsuperscript{1,2}},
%  \textbf{Sixth Author\textsuperscript{1}},
%  \textbf{Seventh Author\textsuperscript{1}},
%  \textbf{Eighth Author \textsuperscript{1,2,3,4}},
%\\
%  \textbf{Ninth Author\textsuperscript{1}},
%  \textbf{Tenth Author\textsuperscript{1}},
%  \textbf{Eleventh E. Author\textsuperscript{1,2,3,4,5}},
%  \textbf{Twelfth Author\textsuperscript{1}},
%\\
%  \textbf{Thirteenth Author\textsuperscript{3}},
%  \textbf{Fourteenth F. Author\textsuperscript{2,4}},
%  \textbf{Fifteenth Author\textsuperscript{1}},
%  \textbf{Sixteenth Author\textsuperscript{1}},
%\\
%  \textbf{Seventeenth S. Author\textsuperscript{4,5}},
%  \textbf{Eighteenth Author\textsuperscript{3,4}},
%  \textbf{Nineteenth N. Author\textsuperscript{2,5}},
%  \textbf{Twentieth Author\textsuperscript{1}}
%\\
%\\
%  \textsuperscript{1}Affiliation 1,
%  \textsuperscript{2}Affiliation 2,
%  \textsuperscript{3}Affiliation 3,
%  \textsuperscript{4}Affiliation 4,
%  \textsuperscript{5}Affiliation 5
%\\
%  \small{
%    \textbf{Correspondence:} \href{mailto:email@domain}{email@domain}
%  }
%}

\begin{document}
\maketitle
% \begin{abstract}
% \end{abstract}

\section{Introduction}

A frequent problem that arises within services that serve large amounts of content to the public is the generation of searchable tags for this content. Some services rely on the creators of the content to create these tags themselves, or may choose to auto-generate these tags, which can be error-prone. It is important that these tags are fairly accurate, as it is often through these tags that users of the service find suitable content to view, either through recommendations or manual searches.

A specific example of this problem is the classification of movies into their genre(s), a problem likely faced by streaming platforms. Should these providers choose to auto-generate these tags, the most efficient method involves analysis of the provided movie description. For the process to be automatic, the tag generation algorithm must include some level of natural language processing, which is where the challenge of the task originates. Movie descriptions are often short, vague, and potentially misleading without greater context, despite often being obvious to human readers as to what genre the movie might be. As for the decision of what genre the movie itself belongs to, people generally can agree, although difficulty can result if a movie belongs to multiple genres at once.

We aim to create a natural language processing algorithm that would take as input the name and description of a movie, and output its genre to aid in this issue.

\section{Related Work}

When searching for a dataset to approach this problem, Arevalo et al.'s "Gated Multimodal Units for Information Fusion" was extremely helpful. It redirected the team towards MovieLens, a database containing millions of movies and their details. The paper presents the Gated Multimodal Unit (GMU), a neural network that learns how much to rely on each input modality through a gating mechanism. The model performs multilabel (multi-genre) classification and combines text from plot summaries and images from movie posters, outperforming both single-modality baselines and other multimodal fusion approaches. (Arevalo et al., 2017)

Do et.al work on the similar task of book genre classification, taking as input the book cover as an image as well as its title. They employ a hybrid model, utilizing a convolutional neural network to process the image, as well as a BERT model to analyze book titles. Combining these features with a dot-product attention mechanism, they were able to achieve impressive performance with an F-score of 0.79. While we do not consider image input within our own work, this work does point towards the importance of textual features in applications such as genre classification as their work improved upon the performance of image and textual features alone. (Do et al., 2020)

This work proposes a Genre Attention Model that would use transformer-based models and hierarchal mechanism to classify movie genres from plot descriptions, a multi-label task. It found that a DeBERTa-based model performed the best and was the most efficient. The study also showed that text-based approaches to a task like this are typically more effective than audio or visual methods. (Zhu et al., 2022)



\section{Dataset}
\subfile{Nour/2-dataset.tex}

\section{Features}

\subfile{Nour/3-LR-features-and-inputs.tex}
\subfile{Dan/3-ZS-features-and-inputs.tex}
\subfile{Lily/3-BBU-features-and-inputs.tex}

\subfile{Sriya-report-sections/3-rnn-features-inputs.tex}

\section{Implementation}

\subfile{Nour/4-LR-implementation.tex}
\subfile{Dan/4-ZS-implementation.tex}
\subfile{Lily/4-BBU-implementation.tex}

\subfile{Sriya-report-sections/4-rnn-implementation.tex}

\section{Results and Evaluation}

\subsection{Evaluation}
\subfile{Nour/5p1-LR-evaluation.tex}

\subsection{Baseline Results}
\subfile{Nour/5p2-baseline-results.tex}

\subsection{Compare Models}
\subsubsection{Participant Models}
One of the approaches taken by a participant is finetuning a DeBERTa (deberta-vs-base) model. 
They finetuned on 90\% of the given data and used the remaining 10\% for validation across 10 epochs. 
They achieved an accuracy of 0.62 and a (weighted) f1 score of 0.59 on the validation set. 
On the other hand, when testing the model on the provided test set, they achieved a (weighted) f1 score of 0.51 and an accuracy of 0.53. 

The second approach taken by another participant is finetuning a RoBERTa (Roberta-base) model. 
As can be seen below, the model scored a f1-score of 0.53 and an accuracy of 0.55. 

% Add model comparison subfiles here
\subfile{Nour/6-LR-compare.tex}
\subfile{Dan/6-ZS-compare.tex}
\subfile{Lily/6-BBU-compare.tex}

\subfile{Sriya-report-sections/6-rnn-comparision.tex}

% Add your model f1 score and accuracy here
\begin{table*}
  \centering
  \begin{tabular}{lll}
    \hline
    \textbf{Model} & \textbf{F1 Score (Weighted)} & \textbf{Accuracy} \\
    \hline
    Majority Baseline & 0.11 & 0.27 \\
    Random Baseline & 0.07 & 0.05 \\
    \hline
    Base Logistic Regression (Word2Vec) & 0.45 & 0.48 \\
    Variation 1: Logistic Regression (TF-IDF) & 0.39 & 0.45 \\
    Variation 2: Logistic Regression (DistilBERT) & 0.51 & 0.53 \\
    \hline
    Zero-Shot (Flan-T5-Base) & 0.37 & 0.39 \\
    \hline
    Fine-Tuned Transformer (BERT-base-uncased) & 0.53 & 0.56 \\
    \hline
    Participant 1: DeBERTa (deberta-vs-base) & 0.51 & 0.53 \\
    Participant 2: RoBERTa (Roberta-base) & 0.53 & 0.55 \\
    \hline
  \end{tabular}
  \caption{\label{model-results}
    Model performance comparison on movie genre classification.
  }
\end{table*}


\section{Error Analysis}

\subfile{Nour/7p1-LR-error-analysis.tex}
\subfile{Dan/7p1-ZS-error-analysis.tex}
\subfile{Lily/7p1-BBU-error-analysis.tex}

\subfile{Sriya-report-sections/7p1-rnn-error-analysis.tex}

\section{Conclusions}

\subfile{Nour/7p2-LR-conclusions.tex}

\subfile{Dan/7p2-ZS-conclusions.tex}

\subfile{Lily/7p2-BBU-conclusions.tex}

\subfile{Sriya-report-sections/7p2-rnn-conclusion.tex}

\nocite{*}
\bibliography{custom}

\section*{Team Contributions}

\subfile{Nour/8-Nour-contributions.tex}

\subfile{Dan/8-Dan-Contributions.tex}

\subfile{Lily/8-Lily-contributions.tex}

\subfile{Sriya-report-sections/8-Sriya-contributions.tex}


\end{document}
